\paragrafos{Επίλυση εξίσωσης της μορφής αx+β=γ}{Α.6.2, «Επίλυση εξίσωσης της μορφής $\alpha x+\beta=\gamma$»}

\begin{anagke}
Η θεία Μαρία δείχνει στην Ζωή τον τελευταίο λογαριασμό του κινητού της: $15$ ευρώ. «Το πακέτο μου έχει πάγιο $9$ ευρώ τον μήνα και κάθε λεπτό ομιλίας εκτός πακέτου κοστίζει $0{,}03$ ευρώ», λέει. «Πόσα λεπτά μίλησα;». Η Ζωή δοκιμάζει: «Εκατό λεπτά, τριάντα εκατό... όχι, τρία ευρώ... και πάλι λάθος». Ο Ηλίας γράφει: αν $x$ είναι τα λεπτά, ο λογαριασμός είναι $9+0{,}03x$, και αυτός ισούται με $15$.

Η Φανή ρωτάει: «Η $9+0{,}03x=15$ δεν είναι σαν τις εξισώσεις της ζυγαριάς. Έχει δύο πράξεις γύρω από το $x$». Η κυρία Αλεξάνδρα απαντά με ερώτηση: «Ποια πράξη έγινε στο $x$ πρώτη και ποια δεύτερη; Και αν γυρίσουμε τη ζυγαριά πίσω, με ποια σειρά θα τις αναιρέσουμε;».

Σε αυτή την παράγραφο θα μάθουμε να λύνουμε εξισώσεις της μορφής $\alpha x+\beta=\gamma$ με τις ιδιότητες της ισότητας, να ελέγχουμε τη λύση με επαλήθευση και να μεταφράζουμε προβλήματα της ζωής σε τέτοιες εξισώσεις.
\end{anagke}

\begin{lexi}{επίλυση, ισοδύναμες, πρόβλημα}{επίλυση, πρόβλημα}
Η \lexiword{επίλυση} είναι η διαδικασία με την οποία «λύνουμε» τον κόμπο, δηλαδή βρίσκουμε τη λύση. Το \lexiword{πρόβλημα} έρχεται από το «προ-βάλλω», δηλαδή «βάζω μπροστά σου» ένα εμπόδιο. Οι εξισώσεις που προκύπτουν σε κάθε βήμα της επίλυσης λέγονται \lexiword{ισοδύναμες}, γιατί έχουν την ίδια λύση.
\end{lexi}

\begin{metafrasi}
Ο λογαριασμός είναι $15$ ευρώ. & $9+0{,}03x=15$\\
Αφαιρώ και από τα δύο μέλη το $9$. & $0{,}03x=6$\\
Διαιρώ και τα δύο μέλη με το $0{,}03$. & $x=200$\\
Το διπλάσιο ενός αριθμού αυξημένο κατά $3$ κάνει $15$. & $2x+3=15$\\
Πέντε ίσα μέρη και $10$ ακόμη κάνουν $70$. & $5x+10=70$\\
\end{metafrasi}

\tmima{Η θεωρία}

\tmimasmall{Η μέθοδος}

Στην εξίσωση $\alpha x+\beta=\gamma$ ο άγνωστος πολλαπλασιάστηκε με το $\alpha$ και μετά προστέθηκε το $\beta$. Για να τον ελευθερώσουμε αναιρούμε τις πράξεις με την αντίστροφη σειρά, με τις ιδιότητες της ισότητας:
\begin{gather*}
\alpha x+\beta=\gamma\\
\alpha x+\beta-\beta=\gamma-\beta\qquad(\text{αφαιρώ και από τα δύο μέλη το }\beta)\\
\alpha x=\gamma-\beta\qquad(\text{κάνω τις πράξεις})\\
\frac{\alpha x}{\alpha}=\frac{\gamma-\beta}{\alpha}\qquad(\text{διαιρώ και τα δύο μέλη με το }\alpha,\ \alpha\neq0)\\
x=\frac{\gamma-\beta}{\alpha}.
\end{gather*}
Όλες οι εξισώσεις που γράψαμε είναι \textbf{ισοδύναμες}: έχουν την ίδια λύση. Η τελευταία δείχνει καθαρά τη λύση.

\begin{center}
\binzyg{\binsx\ \binsx\ \binsw{6}}{\binsw{14}}
\quad$\Longrightarrow$\quad
\binzyg{\binsx\ \binsx}{\binsw{8}}
\quad$\Longrightarrow$\quad
\binzyg{\binsx}{\binsw{4}}
\end{center}
Η ζυγαριά δείχνει την εξίσωση $2x+6=14$. Βγάζω $6$ από κάθε δίσκο και μένει $2x=8$. Χωρίζω κάθε δίσκο σε δύο ίσα μέρη και μένει $x=4$.

\begin{kanonas}{λύση της $\alpha x+\beta=\gamma$}
1. Αφαιρώ και από τα δύο μέλη το $\beta$ (ή προσθέτω, αν το $\beta$ είναι αρνητικό).\\
2. Κάνω τις πράξεις.\\
3. Διαιρώ και τα δύο μέλη με τον συντελεστή $\alpha$ του αγνώστου.\\
4. \textbf{Επαληθεύω} στην αρχική εξίσωση.
\end{kanonas}

\tmimasmall{Όταν υπάρχουν παρενθέσεις ή όμοιοι όροι}

Αν στην εξίσωση υπάρχουν παρενθέσεις, κάνουμε πρώτα την επιμεριστική ιδιότητα και την αναγωγή ομοίων όρων σε κάθε μέλος (Α5), και μετά την αντιμετωπίζουμε ως $\alpha x+\beta=\gamma$. Για παράδειγμα η $3(x-2)+2(x+1)=6$ γίνεται $3x-6+2x+2=6$, δηλαδή $5x-4=6$.

\tmimasmall{Όταν ο άγνωστος χάνεται}

Αν ο συντελεστής $\alpha$ γίνει $0$, δεν μπορούμε να διαιρέσουμε. Η εξίσωση $0\cdot x=\delta$ είναι τότε αδύνατη (αν $\delta\neq0$) ή ταυτότητα (αν $\delta=0$), όπως είδαμε στην προηγούμενη παράγραφο.

\tmimasmall{Προβλήματα με εξίσωση}

\begin{enumerate}[leftmargin=1.6em,itemsep=0.3ex,topsep=2pt]
\item Διαβάζω και βρίσκω τι ζητάω. Τον ονομάζω με ένα γράμμα και γράφω τι σημαίνει, με τις μονάδες του: «Έστω $x$ η τιμή ενός εισιτηρίου, σε ευρώ».
\item Εκφράζω τα άλλα στοιχεία με τη βοήθεια του αγνώστου. Εδώ βοηθάει ένα σχέδιο, π.χ. ένα ορθογώνιο \emph{μοντέλο ράβδου} (δες το παράδειγμα με το ταμείο της εκδρομής).
\item Γράφω την εξίσωση, δηλαδή τη φράση «το ένα ποσό ισούται με το άλλο».
\item Λύνω την εξίσωση.
\item Επαληθεύω στο \emph{πρόβλημα}, όχι μόνο στην εξίσωση, και απαντώ με λόγια και μονάδες.
\end{enumerate}

\begin{prosoxi}
\begin{itemize}[leftmargin=1.2em,itemsep=0.3ex,topsep=2pt]
\item Στην $2x+3=11$ αφαιρώ πρώτα το $3$ και μετά διαιρώ με το $2$: $x=(11-3):2=4$. Δεν γράφω $x=11+3=14$ και δεν διαιρώ το $11$ με το $2$ πριν από την αφαίρεση.
\item Όταν διαιρώ με τον συντελεστή, τον διαιρώ και με το πρόσημό του: από την $-3x=12$ βγαίνει $x=-4$, όχι $4$.
\item Μια λύση που δεν βγάζει νόημα στο πρόβλημα (π.χ. $-5$ μαθητές ή $2{,}5$ παιδιά) σημαίνει ότι κάτι δεν πήγε καλά.
\end{itemize}
\end{prosoxi}

\tmima{Λυμένα παραδείγματα}

\paradeigma{Δύο βήματα}
\tixero Η εξίσωση $2x+6=14$.\par
\tithelo Να τη λύσω και να επαληθεύσω.\par
\possindeo $2x+6=14$, άρα $2x+6-6=14-6$, δηλαδή $2x=8$. Διαιρώ με το $2$: $x=4$.\par
\apantisi $x=4$.\par
\elegxos $2\cdot4+6=14$. Σωστό.

\paradeigma{Αρνητικά, δεκαδικά, κλάσματα}
\tixero Οι εξισώσεις $7-3x=-5$, $\dfrac{1}{2}x-\dfrac{3}{4}=\dfrac{1}{4}$ και $1{,}5x-2{,}25=3$.\par
\tithelo Να τις λύσω.\par
\possindeo Πρώτη: $-3x=-5-7=-12$, άρα $x=-12:(-3)=4$. Δεύτερη: $\frac{1}{2}x=\frac{1}{4}+\frac{3}{4}=1$, άρα $x=2$. Τρίτη: $1{,}5x=3+2{,}25=5{,}25$, άρα $x=5{,}25:1{,}5=3{,}5$.\par
\apantisi $4$, $2$ και $3{,}5$.\par
\elegxos $7-12=-5$. $\frac{1}{2}\cdot2-\frac{3}{4}=\frac{1}{4}$. $1{,}5\cdot3{,}5-2{,}25=5{,}25-2{,}25=3$.

\paradeigma{Με παρενθέσεις}
\tixero Η εξίσωση $3(x-2)+2(x+1)=6$.\par
\tithelo Να τη λύσω.\par
\possindeo Επιμεριστική ιδιότητα: $3x-6+2x+2=6$. Αναγωγή ομοίων όρων: $5x-4=6$. Προσθέτω $4$: $5x=10$. Διαιρώ με $5$: $x=2$.\par
\apantisi $x=2$.\par
\elegxos $3\cdot0+2\cdot3=6$. Σωστό.

\paradeigma{Ο άγνωστος χάνεται}
\tixero Οι εξισώσεις $2(x+3)=2x+5$ και $3(x+1)-3x=3$.\par
\tithelo Να τις λύσω.\par
\possindeo Η πρώτη γίνεται $2x+6=2x+5$ και, αφαιρώντας $2x$ και από τα δύο μέλη, $6=5$, που δεν ισχύει ποτέ. Η δεύτερη γίνεται $3x+3-3x=3$, δηλαδή $3=3$, που ισχύει πάντα.\par
\apantisi Η πρώτη είναι αδύνατη και η δεύτερη ταυτότητα.\par
\elegxos Με $x=1$ η πρώτη δίνει $8\neq7$ και η δεύτερη $3=3$. Με $x=5$ η πρώτη $16\neq15$ και η δεύτερη $3=3$.

\paradeigma{Ο λογαριασμός του κινητού}
\tixero Το πακέτο της θείας Μαρίας κοστίζει $9$ € πάγιο και $0{,}03$ € το λεπτό. Ο λογαριασμός ήταν $15$ €.\par
\tithelo Πόσα λεπτά μίλησε.\par
\possindeo Έστω $x$ τα λεπτά. Ο λογαριασμός είναι $9+0{,}03x=15$. Άρα $0{,}03x=6$ και $x=6:0{,}03=200$.\par
\apantisi $200$ λεπτά.\par
\elegxos $9+0{,}03\cdot200=9+6=15$ €. Σωστό.

\paradeigma{Το ταμείο της εκδρομής}
\tixero Τα $24$ παιδιά της Α3 πληρώνουν από $x$ ευρώ το καθένα και το σχολείο δίνει ακόμη $60$ €. Το σύνολο είναι $420$ €.\par
\tithelo Πόσο πληρώνει κάθε παιδί.\par
\possindeo Το μοντέλο ράβδου:
\begin{center}
\begin{tikzpicture}[x=0.19cm]
\foreach \i in {0,...,23}{\draw[fill=gblue!25] (\i*1.0,0) rectangle (\i*1.0+1,0.5);}
\draw[fill=black!15] (24,0) rectangle (30,0.5);
\node at (27,0.25) {\scriptsize 60};
\draw[<->] (0,-0.35)--(24,-0.35) node[midway,below,font=\small] {$24$ φορές το $x$};
\draw[<->] (0,0.95)--(30,0.95) node[midway,above,font=\small] {$420$};
\end{tikzpicture}
\end{center}
Άρα $24x+60=420$, $24x=360$ και $x=15$.\par
\apantisi $15$ € το παιδί.\par
\elegxos $24\cdot15=360$ και $360+60=420$. Σωστό.

\paradeigma{Οι σελίδες του βιβλίου}
\tixero Η Ρήνη διάβασε τα $\frac{2}{5}$ ενός βιβλίου και άλλες $30$ σελίδες. Έχει διαβάσει συνολικά $150$ σελίδες.\par
\tithelo Πόσες σελίδες έχει το βιβλίο.\par
\possindeo Έστω $x$ οι σελίδες. Τότε $\frac{2}{5}x+30=150$. Άρα $\frac{2}{5}x=120$ και $x=120\cdot\frac{5}{2}=300$.\par
\apantisi $300$ σελίδες.\par
\elegxos Τα $\frac{2}{5}$ των $300$ είναι $120$ και $120+30=150$.

\paradeigma{Η ηλικία της γιαγιάς}
\tixero Η γιαγιά Μυρτώ είναι $70$ ετών, δέκα χρόνια περισσότερα από το πενταπλάσιο της ηλικίας ενός εγγονού της.\par
\tithelo Την ηλικία του εγγονού.\par
\possindeo Έστω $x$ η ηλικία του εγγονού. Τότε $5x+10=70$, $5x=60$ και $x=12$.\par
\apantisi $12$ ετών.\par
\elegxos $5\cdot12+10=70$.

\tmima{Ασκήσεις}

\zoni{Το κατάλαβα}{τα βήματα και η επαλήθευση}

\ex Λύσε τις εξισώσεις: α) $3x+4=19$ \quad β) $6x+8=62$ \quad γ) $5+12\alpha=77$ \quad δ) $9\alpha-15=111$.

\ex Λύσε τις εξισώσεις: α) $-3x-5=10$ \quad β) $2-5x=-13$ \quad γ) $12x+31=7$ \quad δ) $8+6x-2=7-7$.

\ex Βάλε στη σωστή σειρά τα βήματα επίλυσης της $3x-7=11$: α) Διαιρώ και τα δύο μέλη με το $3$. \quad β) Προσθέτω $7$ και στα δύο μέλη. \quad γ) Επαληθεύω: $3\cdot6-7=11$. \quad δ) Έχω $3x=18$. \quad ε) Βρίσκω $x=6$.

\ex Συμπλήρωσε τη λύση: $5x+10=45$, άρα $5x=45-\dots=\dots$ και $x=\dots:5=\dots$.

\ex Εξέτασε αν ο αριθμός είναι λύση της εξίσωσης: α) $4x-5=7$, $x=3$ \quad β) $3x+8=2$, $x=-2$ \quad γ) $2x-1=9$, $x=4$ \quad δ) $4x+3=5$, $x=\dfrac{1}{2}$.

\ex Αξιοποίησε την επιμεριστική ιδιότητα και λύσε: α) $2(x-1)=18$ \quad β) $3-(x-2)=7$ \quad γ) $4(x+1)-2(x-3)=20$ \quad δ) $x-2(x-4)=10$.

\ex Λύσε: α) $\dfrac{1}{2}x+3=7$ \quad β) $\dfrac{2}{3}x-1=5$ \quad γ) $0{,}5x+1{,}2=4$ \quad δ) $\dfrac{x}{8}-1=0$.

\zoni{Το χρησιμοποιώ}{ιστορίες που γίνονται εξισώσεις}

\ex Γράψε εξίσωση και λύσε: α) το διπλάσιο ενός αριθμού αυξημένο κατά $3$ κάνει $15$ \quad β) το τριπλάσιο ενός αριθμού ελαττωμένο κατά $7$ κάνει $20$ \quad γ) το $56$ ελαττωμένο κατά το τετραπλάσιο ενός αριθμού κάνει $8$ \quad δ) το μισό ενός αριθμού αυξημένο κατά $3$ κάνει $10$.

\ex \textbf{Το ποδήλατο.} Ένα ποδήλατο κοστίζει $180$ €. Ο Ηλίας έχει $45$ € και αποταμιεύει $15$ € την εβδομάδα. Σε πόσες εβδομάδες θα το αγοράσει;

\ex \textbf{Η εκδρομή.} Τα $25$ παιδιά μιας τάξης πληρώνουν από $x$ ευρώ και το σχολείο δίνει ακόμη $50$ €. Το σύνολο είναι $300$ €. Σχεδίασε το μοντέλο ράβδου, γράψε την εξίσωση και βρες το $x$.

\ex \textbf{Τα καφάσια.} Ο παππούς Γιάννης έχει δύο ίδια καφάσια με ελιές και ακόμη $3$ κιλά ελιές ξεχωριστά. Συνολικά έχει $31$ κιλά. Πόσα κιλά έχει κάθε καφάσι;

\ex \textbf{Οι ηλικίες.} Η θεία Μαρία είναι $38$ ετών, δύο χρόνια περισσότερα από το τριπλάσιο της ηλικίας του Νίκου. Πόσο ετών είναι ο Νίκος;

\zoni{Το σκέφτομαι}{ανάποδη πορεία, λάθη, δικά σου προβλήματα}

\ex Ο Στέλιος λύνει την $2x+3=11$ γράφοντας $x=11+3=14$ και μετά $x=14:2=7$. Βρες το λάθος, γράψε τη σωστή λύση και επαλήθευσε.

\ex Σκέφτηκα έναν αριθμό, τον πολλαπλασίασα επί $3$, αφαίρεσα $4$, διαίρεσα με το $2$ και βρήκα $10$. Βρες τον αριθμό πηγαίνοντας ανάποδα και μετά γράψε την εξίσωση.

\ex Γράψε ένα πρόβλημα της καθημερινής ζωής που λύνεται με την εξίσωση $2x+4=20$ και λύσε το.

\ex \textbf{Η μυστική συνταγή της γιαγιάς.} Η γιαγιά Μυρτώ φτιάχνει κουλουράκια με ζάχαρη, βούτυρο και αμύγδαλα. Το βούτυρο είναι διπλάσιο από τη ζάχαρη και τα αμύγδαλα είναι $10$ g λιγότερα από τη ζάχαρη. Όλα τα υλικά ζυγίζουν $390$ g. Πόσα γραμμάρια είναι το καθένα;

\zoni{Για το σπίτι}{φτιάξε και λύσε}

\ex Σκέψου ένα πρόβλημα από το σπίτι σου που λύνεται με την εξίσωση $4x+5=45$ και λύσε το.

\ex Λύσε την εξίσωση $2x-3(0{,}5x+1)-x=1$ και κάνε την επαλήθευση.
